\documentclass[11pt]{article}
\usepackage{mathblatt}
% vorspann.tex – Kompetenzblatt, zusammenbau v0.7 (2026-09-28).
% Ergänzt mathblatt.sty für Kompetenzblätter, ohne sie zu ändern
% (layout-befunde.md 1–34 vom 28.09.). Wird nach \usepackage{mathblatt}
% eingelesen.
\makeatletter
\RequirePackage{newunicodechar}
% Zeichen, die Latin Modern nicht hat (Befund 20): Text aus einer
% Ersatzschrift, Mathe als Befehl; ohne Ersatzschrift auch Text als Mathe.
\newif\ifkbersatz
\IfFontExistsTF{FreeSerif}{\newfontfamily\kbersatzschrift{FreeSerif}\kbersatztrue}{%
 \IfFontExistsTF{DejaVu Serif}{\newfontfamily\kbersatzschrift{DejaVu Serif}\kbersatztrue}{%
  \IfFontExistsTF{Cambria Math}{\newfontfamily\kbersatzschrift{Cambria Math}\kbersatztrue}{%
   \IfFontExistsTF{Segoe UI Symbol}{\newfontfamily\kbersatzschrift{Segoe UI Symbol}\kbersatztrue}{}}}}
\newcommand{\kbzeichen}[2]{\ifmmode #2\else\ifkbersatz{\kbersatzschrift #1}\else\ensuremath{#2}\fi\fi}
\newcommand{\kbzeichenhoch}[1]{\ifkbersatz\ifmmode\text{\kbersatzschrift #1}\else{\kbersatzschrift #1}\fi\else ?\fi}
% Kopf- und Fußzeile (Befund 1, 2): Kopf Thema · Blattart, Fuß Kennung und Seite
\newcommand{\kbkopfzeile}[2]{\pagestyle{fancy}\fancyhf{}\renewcommand{\headrulewidth}{0pt}%
  \fancyhead[L]{\footnotesize\color{mbgrau}#1}%
  \fancyfoot[L]{\footnotesize\color{mbgrau}#2}%
  \fancyfoot[R]{\footnotesize\color{mbgrau}Seite \thepage}}
% Flattersatz überall (Befund 25)
\raggedright
% Titel des Blatts: Ich-kann-Satz, darunter Zweigzeile
\newcommand{\kbtitel}[2]{\par\noindent{\Large\bfseries\raggedright #1\par}%
  \ifblank{#2}{}{\par\smallskip\noindent{\small #2}\par}\medskip}
% Abschnitt: „Das kennst du schon“, „Zum Merken“
\newcommand{\kbabschnitt}[1]{\par\addvspace{12pt}\Needspace*{5\baselineskip}%
  \noindent{\bfseries #1}\par\nobreak\vspace{1pt}\noindent{\color{mbgrau}\rule{\linewidth}{0.4pt}}\par\nobreak}
% Hauptnummer: Titel hängt am ersten Block; jeder Block (Teilaufgabe) bricht
% nicht, passt er nicht mehr, rückt er ganz auf die nächste Seite (Befund 21).
\newif\ifkbtitel
\newsavebox{\kbbox}
\newlength{\kbhoehe}
\newlength{\kbnummerabstand}\setlength{\kbnummerabstand}{10pt}
\newlength{\kbteilabstand}\setlength{\kbteilabstand}{6pt}
\newenvironment{kbaufgabe}[1]{\stepcounter{aufgabe}\setcounter{teil}{0}%
  \gdef\kbtiteltext{\noindent\hangindent1.6em\hangafter1\makebox[1.6em][l]{\textbf{\theaufgabe.}}#1\par}%
  \global\kbtiteltrue\par\addvspace{\kbnummerabstand}}{\par}
\newenvironment{kbblock}{\par\addvspace{\kbteilabstand}\begin{lrbox}{\kbbox}%
  \begin{minipage}[t]{\linewidth}\raggedright\parskip\z@
  \ifkbtitel\kbtiteltext\vspace{2pt}\global\kbtitelfalse\fi
  \list{}{\leftmargin1.6em\labelwidth1.4em\labelsep0.2em\topsep\z@\partopsep\z@
    \parsep\z@\itemsep\z@\rightmargin\z@}\raggedright}%
  {\endlist\end{minipage}\end{lrbox}%
  \setlength{\kbhoehe}{\dimexpr\ht\kbbox+\dp\kbbox\relax}%
  \Needspace*{\kbhoehe}\noindent\usebox{\kbbox}\par}
% Paarsatz (Platz nutzen, Befund 27): zwei Hauptnummern oder zwei
% Teilaufgaben mit Zeichenaufgabe nebeneinander, als ein Block
\newcommand{\kbnummer}[1]{\stepcounter{aufgabe}\setcounter{teil}{0}%
  \noindent\hangindent1.6em\hangafter1\makebox[1.6em][l]{\textbf{\theaufgabe.}}#1\par\vspace{2pt}}
\newcommand{\kbhalb}[1]{\list{}{\leftmargin1.6em\labelwidth1.4em\labelsep0.2em\topsep\z@
  \partopsep\z@\parsep\z@\itemsep\z@}\raggedright #1\endlist}
\newcommand{\kbpaar}[2]{\par\addvspace{\kbnummerabstand}\begin{lrbox}{\kbbox}%
  \begin{minipage}[t]{0.485\linewidth}\raggedright\parskip\z@ #1\end{minipage}\hfill
  \begin{minipage}[t]{0.485\linewidth}\raggedright\parskip\z@ #2\end{minipage}\end{lrbox}%
  \setlength{\kbhoehe}{\dimexpr\ht\kbbox+\dp\kbbox\relax}\Needspace*{\kbhoehe}\noindent\usebox{\kbbox}\par}
\newcommand{\kbteilpaar}[2]{\item[]\noindent
  \begin{minipage}[t]{0.485\linewidth}\kbhalb{#1}\end{minipage}\hfill
  \begin{minipage}[t]{0.485\linewidth}\kbhalb{#2}\end{minipage}\par}
% Anweisung der Hauptnummer, eigene Zeile über der Teilaufgabe (Befund 5)
\newcommand{\kbanweisung}[1]{\item[]#1\par\vspace{2pt}}
% Kopfzeile einer Teilaufgabe: Buchstabe, Auftakt (halbfett oder Kontext),
% Prüfkennung klein rechts (Befund 3, 12, Beschluss c)
\newlength{\kbkennbreite}\setlength{\kbkennbreite}{1.9cm}
\newcommand{\kbkennung}[1]{{\footnotesize #1}}
\newcommand{\kbteil}[3]{\item[#1]\ifblank{#3}%
  {\parbox[t]{\linewidth}{\raggedright #2}}%
  {\parbox[t]{\dimexpr\linewidth-\kbkennbreite\relax}{\raggedright #2}\hfill
   \parbox[t]{\kbkennbreite}{\raggedleft\kbkennung{#3}}}\par}
% Frage in eigener Zeile (Befund 4), ein Satz je Zeile (Befund 10)
\newcommand{\kbfrage}[1]{\par\vspace{2pt}#1\par}
% Antwortfeld in eigener Zeile darunter, linksbündig (Befund 4, 8)
\newcommand{\kbantwort}[1]{\par\vspace{4pt}\noindent#1\par}
% Zwei Spalten: links Grafik/Tabelle/Aufgabe, rechts Antwort oder Rechenraum
% (Befund 27); oben bündig. Beginnt einen eigenen Listenpunkt ohne Marke.
\newcommand{\kbzweispaltig}[3][0.48]{\item[]\noindent
  \begin{minipage}[t]{#1\linewidth}\raggedright #2\end{minipage}\hfill
  \begin{minipage}[t]{\dimexpr0.98\linewidth-#1\linewidth\relax}\raggedright #3\end{minipage}\par}
% Linke Spalte mit Teilaufgabe (Buchstabe hängt links heraus wie in der Liste)
\newcommand{\kbspalte}[1]{\list{}{\leftmargin\z@\labelwidth1.4em\labelsep0.2em\topsep\z@
  \partopsep\z@\parsep\z@\itemsep\z@}\raggedright #1\endlist}
% Antwortfelder: 2,2 cm, in Punkten 1,2 cm; ohne Umbruch davor
\newcommand{\kbfeld}[1][]{\mbox{\underline{\hspace{2.2cm}}\ifblank{#1}{}{\,#1}}}
\newcommand{\kbfeldk}[1][]{\mbox{\underline{\hspace{1.2cm}}\ifblank{#1}{}{\,#1}}}
% Grafik oder Tabelle unter dem Text, linksbündig (Befund 8, 28)
\newcommand{\kbdarunter}[1]{\par\vspace{4pt}\noindent#1\par}
% Ankreuzen: je Option eine Zeile, linksbündig (Befund 6, 7)
\newcommand{\kbkreuz}[1]{\par\vspace{2pt}\noindent$\square$\ #1\par}
% Bearbeitungsraum (Befund 9): Schreibzeilen wie mathblatt (9 mm)
\newcommand{\kbraum}[1]{\schreibzeilen{#1}}
% Wertetabelle mit Köpfen im Textmodus (Befund 26): wie \wertetabelle der
% Vorlage, aber die Köpfe setzt der Aufruf selbst ($x$, „Zeit in h“).
\NewDocumentCommand{\kbwertetabelle}{o m m m}{%
  \def\mbkopf{}\def\mbwerte{}\def\mbleer{}\setcounter{mbn}{0}%
  \foreach \v in {#4}{\xappto\mbkopf{& $\v$}\xappto\mbleer{& }\stepcounter{mbn}\xdef\mbnn{\arabic{mbn}}}%
  \IfNoValueTF{#1}{}{\foreach \v in {#1}{\ifdefempty{\v}{\xappto\mbwerte{& }}{\xappto\mbwerte{& $\v$}}}}%
  \begingroup\renewcommand{\arraystretch}{1.3}%
  \edef\mbpre{|>{\noexpand\columncolor{mbkasten}}l|*{\mbnn}{>{\noexpand\centering\noexpand\arraybackslash}p{\noexpand\mbzell}|}}%
  \expandafter\mbtabstart\expandafter{\mbpre}\hline
  #2 \mbkopf \\ \hline
  \IfNoValueTF{#1}{\rule{0pt}{\mbzeile}#3 \mbleer \\ \hline}{#3 \mbwerte \\ \hline}
  \end{tabular}\endgroup}
% Merkkasten am Ende (Aufbau des Kompetenzblatts)
\newcommand{\kbmerk}[2]{\par\addvspace{14pt}\begin{lrbox}{\kbbox}\begin{minipage}[t]{\linewidth}%
  \noindent{\bfseries #1}\par\vspace{1pt}\noindent{\color{mbgrau}\rule{\linewidth}{0.4pt}}\par\vspace{4pt}%
  \uebersichtskasten{\raggedright #2}\end{minipage}\end{lrbox}%
  \setlength{\kbhoehe}{\dimexpr\ht\kbbox+\dp\kbbox\relax}\Needspace*{\kbhoehe}\noindent\usebox{\kbbox}\par}
% Lösungsblatt: Nummer und Buchstabe halbfett, Lösung daneben
\newcommand{\kbloesung}[2]{\par\addvspace{3pt}\noindent\hangindent2.8em\hangafter1\makebox[2.8em][l]{\textbf{#1}}#2\par}
% Zeichen dieses Blatts, die Latin Modern nicht hat
\newunicodechar{α}{\kbzeichen{α}{\alpha}}
\newunicodechar{β}{\kbzeichen{β}{\beta}}
\newunicodechar{γ}{\kbzeichen{γ}{\gamma}}
\newunicodechar{≈}{\kbzeichen{≈}{\approx}}
\makeatother

\begin{document}
\kbkopfzeile{Trigonometrie · Kompetenzblatt · Lösungen}{TRI-K1}
% Lösungen TRI-K1
\kbtitel{Lösungen}{Ich kann Seiten und Winkel mit dem Sinussatz berechnen.}
\kbloesung{1.}{$180^\circ - 45^\circ - 85^\circ = 50^\circ$}
\kbloesung{2.}{$x = 3 \cdot 5 = 15$}
\kbloesung{3.}{$3\,000$ m}
\kbloesung{4.}{sin, cos, tan}
\kbloesung{5.}{$\frac{b}{\mathrm{sin}\,51^\circ} = \frac{7}{\mathrm{sin}\,68^\circ}$, $b = 7 \cdot \mathrm{sin}\,51^\circ : \mathrm{sin}\,68^\circ \approx 5{,}9$ cm}
\kbloesung{6.}{$PQ$ liegt gegenüber $R$, $PR$ gegenüber $Q$: $PR = 8 \cdot \mathrm{sin}\,53^\circ : \mathrm{sin}\,69^\circ \approx 6{,}8$ cm}
\kbloesung{7.}{$\gamma = 180^\circ - 63^\circ - 48^\circ = 69^\circ$; $c = 8 \cdot \mathrm{sin}\,69^\circ : \mathrm{sin}\,63^\circ \approx 8{,}4$ cm}
\kbloesung{8.}{$c = 9 \cdot \mathrm{sin}\,24^\circ : \mathrm{sin}\,127^\circ \approx 4{,}6$ cm}
\kbloesung{9.}{$\gamma = 180^\circ - 27^\circ - 118^\circ = 35^\circ$; $c = 7 \cdot \mathrm{sin}\,35^\circ : \mathrm{sin}\,27^\circ \approx 8{,}8$ cm}
\kbloesung{10.}{$b = 4{,}35 \cdot \mathrm{sin}\,47{,}6^\circ : \mathrm{sin}\,58{,}2^\circ \approx 3{,}8$ m}
\kbloesung{11.}{$b = 9 \cdot \mathrm{sin}\,44^\circ : \mathrm{sin}\,77^\circ \approx 6{,}42$ cm, also rund $6{,}4$ cm}
\kbloesung{12.}{$\gamma = 71^\circ$; $a = 10 \cdot \mathrm{sin}\,53^\circ : \mathrm{sin}\,71^\circ \approx 8{,}4$ cm, $b = 10 \cdot \mathrm{sin}\,56^\circ : \mathrm{sin}\,71^\circ \approx 8{,}8$ cm – die Aussage ist falsch.}
\kbloesung{13.}{$\angle ACB = 90^\circ - 57^\circ = 33^\circ$, $\angle ACD = 91^\circ - 33^\circ = 58^\circ$; $AD = 8 \cdot \mathrm{sin}\,58^\circ : \mathrm{sin}\,76^\circ \approx 7{,}0$ m}
\kbloesung{14.}{$a = 16 \cdot \mathrm{sin}\,31^\circ : \mathrm{sin}\,104^\circ \approx 8{,}49$ m; $BP = a - 5 \approx 3{,}5$ m}
\kbloesung{15.}{Sinussatz: $b = 8 \cdot \mathrm{sin}\,53^\circ : \mathrm{sin}\,67^\circ \approx 6{,}9$ cm. Höhe: $h = 8 \cdot \mathrm{sin}\,53^\circ \approx 6{,}39$ cm, $b = h : \mathrm{sin}\,67^\circ \approx 6{,}9$ cm – gleich}
\kbloesung{16.}{$\mathrm{sin}\,\beta = 7 \cdot \mathrm{sin}\,68^\circ : 9 \approx 0{,}7211$, $\beta \approx 46{,}1^\circ$}
\kbloesung{17.}{$c^2 = 5^2 + 8^2 - 2 \cdot 5 \cdot 8 \cdot \mathrm{cos}\,47^\circ$, $c \approx 5{,}9$ cm}
\kbloesung{18.}{$\mathrm{cos}\,\gamma = \frac{5^2 + 7^2 - 6^2}{2 \cdot 5 \cdot 7}$, $\gamma \approx 57{,}1^\circ$}
\kbloesung{19.}{$h_c = b \cdot \mathrm{sin}\,\alpha$ und $h_c = a \cdot \mathrm{sin}\,\beta$, also $b \cdot \mathrm{sin}\,\alpha = a \cdot \mathrm{sin}\,\beta$; beide Seiten durch $\mathrm{sin}\,\alpha \cdot \mathrm{sin}\,\beta$ teilen.}
\kbloesung{20a)}{dritter Winkel $180^\circ - 6^\circ - 137^\circ = 37^\circ$ (gegenüber $220$ cm); $y = 220 \cdot \mathrm{sin}\,137^\circ : \mathrm{sin}\,37^\circ \approx 249{,}3$ cm}
\kbloesung{20b)}{$\gamma = 180^\circ - 31^\circ - 113^\circ = 36^\circ$; $BC = 520 \cdot \mathrm{sin}\,31^\circ : \mathrm{sin}\,36^\circ \approx 455{,}6$ m}
\kbloesung{20c)}{$\angle ABD = 42^\circ$, $\angle DBC = 62^\circ$; $BC = 1\,300 \cdot \mathrm{sin}\,57^\circ : \mathrm{sin}\,62^\circ \approx 1\,235$ m – nicht gleich $1\,300$ m, die Aussage ist falsch.}
\kbloesung{20d)}{Winkel bei $D$: $39^\circ$; $DE = 9 \cdot \mathrm{sin}\,118^\circ : \mathrm{sin}\,39^\circ \approx 12{,}63$ m; $DP = DE - 8 \approx 4{,}6$ m}
\kbloesung{20e)}{$\angle ACB = 31^\circ$, $\angle ACD = 57^\circ$; $AD = 7 \cdot \mathrm{sin}\,57^\circ : \mathrm{sin}\,78^\circ \approx 6{,}0$ m}
\end{document}
